% Source-sized Hintikka unit; original1953 and distinct historical reading copy.
\chapter{Hintikka: constituents, prediction, and inductive generalization}
\label{sec:p15-hintikka-inductive-logic}

\section{The problem is the law, not just its next instance}
\label{sec:p15-hintikka-program}

Carnap's instance confirmation and Reichenbach's frequency posit answer
different questions from the probability of a genuine universal law.
Hintikka's program makes that last question explicit: describe which
kinds of individuals exist, assign probabilities to those descriptions,
and update them on a sample. Its distinctive step is to distribute mass
over \emph{constituents}, rather than let an unrestricted continuous
frequency distribution determine the mass of every law.
Niiniluoto's historical account distinguishes the Jerusalem presentation
of 1964, published in 1965; the combined system of 1965; the two-dimensional
continuum of 1966; and the later Hintikka--Niiniluoto axiomatization
\cite[pp.~311--325]{Niiniluoto2011HintikkaProgramReading}.

There are two source roles in this chapter. Hintikka's 1953 original
supplies the distributive-normal-form language machinery. Niiniluoto's
historical author copy supplies the account of the later probability
constructions and their development. Its 2009 footer distinguishes that
copy from the final published 2011 version. The constructions below have
their own declared domain and normalization; the original 1965/1966
presentations and 1976/1980 axiomatization are distinct source presentations.

The historian interprets the program as retaining probability-based
induction and the distinction between epistemic and frequency probability,
but relaxing the claim that inductive probability is independent of
extra-logical factors. Language, contextual regularity assumptions and
inductive parameters may affect the method. Once they are fixed, a
mathematical calculation can be exact without making their choice a
logical truth. This is an argument about the role of the calculus, not
a theorem that one preferred choice follows from logic alone.

\section{Kinds of individuals and constituents}
\label{sec:p15-hintikka-constituents}

\begin{hbdefinition}{def:p15-hintikka-monadic-language}{The nonempty monadic setting}
Take a finite list of unary predicates \(M_1,\ldots,M_k\), no identity
predicate, no function symbols and no relational predicates of higher
arity. The domain is nonempty. Let \(K=2^k\) and let
\(Q_1,\ldots,Q_K\) be all conjunctions choosing \(M_jx\) or
\(\neg M_jx\) for each \(j\). More generally one may begin with
any specified finite exhaustive, mutually exclusive classification into
\(K\) cells. Individual names \(a_1,a_2,\ldots\) are used for
sample reports. A \emph{generalization} here is a closed sentence without
individual names. Without identity its truth cannot distinguish one
occupant of a cell from several occupants.

For \(\varnothing\ne S\subseteq\{1,\ldots,K\}\), define
\[
 C_S\defeq
 \bigwedge_{i\in S}\exists x\,Q_i(x)
 \ \land\ \forall x\bigvee_{i\in S}Q_i(x).
\]
Its \emph{width} is \(w=|S|\). It says exactly which cells are
nonempty, not their population proportions. The maximally wide
\(C_{\{1,\ldots,K\}}\) is the \emph{atomistic constituent}.
The empty-subset expression is false on this declared nonempty domain.
\end{hbdefinition}

The original 1953 discussion uses red and round as its two predicates.
Its four individual profiles are neither, round only, red only and both.
A constituent first asserts the existence of selected profiles and then
excludes all others. The original also permits an empty-domain
alternative in its initial list; the present nonempty-domain convention
is explicit, not silently attributed to that list
\cite[sec.~3, subsecs.~16,20, pp.~28--31,37--39]{Hintikka1953DistributiveNormalForms}.

\begin{hbtheorem}{thm:p15-hintikka-monadic-normal-form}{proved=annex:p15.hintikka-monadic-normal-form}{The monadic normal form and its inferential content}
The \(Q_i\) partition all individual profiles. The \(2^K-1\)
nonempty-domain constituents are consistent, mutually exclusive and
jointly exhaustive. Every closed monadic sentence without identity or
individual names is equivalent to a disjunction of a uniquely determined
subset of them. For a Boolean property \(A(x)\) allowing precisely the
cells in \(B\),
\[
 C_S\models\forall x A(x)\ \Longleftrightarrow\ S\subseteq B,
 \qquad
 C_S\models\exists x A(x)\ \Longleftrightarrow\ S\cap B\ne\varnothing.
\]
If a nonempty sample \(e\) has observed support \(T\), with
\(|T|=c\geq1\), then precisely the \(2^{K-c}\) constituents with
\(T\subseteq S\) are compatible with \(e\). Empty evidence instead
leaves all \(2^K-1\) nonempty-domain constituents compatible. In either
case there are \(\binom{K-c}{w-c}\) compatible constituents of width
\(w\), for \(\max(1,c)\leq w\leq K\).
\end{hbtheorem}

Thus an existential report about a previously unseen cell eliminates
every constituent which excludes it. Repeating a known cell eliminates
none of these surviving constituents but can change their relative
probabilities. These are the eliminative and enumerative contributions
of evidence. The normal form is not by itself a probability assignment.

\section{The original normal-form method beyond unary languages}
\label{sec:p15-hintikka-original-normal-forms}

Hintikka's 1953 original does not stop at the unary setting. It defines
constituents recursively over a finite basis of predicates and a fixed
quantificational degree. Its first-kind forms make independent positive
or negative choices for atomic components and existentially quantified
subordinate constituents; its second-kind forms rewrite the negative
existentials as a universal restriction to the permitted alternatives.
The original's degree is maximum quantifier nesting, not the total
number of quantifiers. The following equivalent readable construction
keeps old atomic components at each level; it is more redundant than
the source's grouping by highest variable index
\cite[sec.~3, subsecs.~17--21, pp.~31--39]{Hintikka1953DistributiveNormalForms}.

\begin{hbdefinition}{def:p15-hintikka-depth-types}{Recursive descriptions at bounded degree}
Fix a finite relational vocabulary without identity or functions and
an integer \(d\geq0\). At level \(a\leq d\), let \(A_a\)
be all atomic formulas over variables \(x_1,\ldots,x_a\), including
repeated arguments and any zero-ary predicates. Let \(\mathcal T_d\)
contain the complete Boolean conjunctions over \(A_d\).
Recursively let \(\mathcal T_a\), for \(a<d\), contain all
complete Boolean conjunctions over the finite list
\[
 A_a\ \cup\ \{\exists x_{a+1}\tau:
                  \tau\in\mathcal T_{a+1}\}.
\]
Both signs are chosen for every listed component. A conjunction in
\(\mathcal T_a\) may be inconsistent. No inconsistency test is
built into this syntactic enumeration.
\end{hbdefinition}

\begin{hbtheorem}{thm:p15-hintikka-depth-normal-form}{proved=annex:p15.hintikka-depth-normal-form}{Bounded-degree relational normal forms}
At each level the descriptions are mutually exclusive and exhaustive
as semantic alternatives, although some may have no model. Every formula
with free variables among \(x_1,\ldots,x_a\) and quantifier nesting
at most \(d-a\) is equivalent to a disjunction of descriptions in
\(\mathcal T_a\). In particular every closed formula of degree at
most \(d\) has such a form over \(\mathcal T_0\).
This is a finite effective syntactic transformation. It is not a
procedure deciding which relational descriptions are satisfiable, nor
a normalized inductive probability on all first-order theories.
\end{hbtheorem}

The distinction matters. A purported probability which simply gives
each syntactically listed relational description positive equal mass
would give positive mass to some contradictions. The original already
warns that removing the redundant parts of its unrestricted relational
forms is tied to the decision problem. The historical chapter reports
Hilpinen's depth-two development and says the general inductive case had
not been studied, rather than claiming that the normal-form theorem
settles that probability construction
\cite[introduction, pp.~7--9]{Hintikka1953DistributiveNormalForms}
\cite[pp.~326--329]{Niiniluoto2011HintikkaProgramReading}.

\section{A declared probability construction on constituents}
\label{sec:p15-hintikka-probability-construction}

The historian describes the Jerusalem system by equal constituent mass,
then equal mass among the states belonging to a constituent; the combined
system inserts an equal-structure step between those two allocations.
The two-dimensional system separates a parameter for singular prediction
from a parameter for generalization. These are different allocations,
not three names for Bayes's formula. Its finite-universe descriptions
and infinite-universe limiting probabilities also require a distinction.
Here is an explicit countable-domain realization of the displayed
constituent-conditioned predictive formulas. Its existence is proved,
rather than inferred from the name of the system.

\begin{hbdefinition}{def:p15-hintikka-conditional-prediction}{Constituent priors and conditional sample law}
Assign each nonempty \(S\) a positive prior \(\pi_S\), with
\(\sum_{S\ne\varnothing}\pi_S=1\). Choose positive parameters
\(\lambda_w\), one for each width. Conditional on \(C_S\),
give an ordered sample word \(e=(i_1,\ldots,i_n)\) likelihood
\[
 L_S(e)\defeq
 \begin{cases}
 \displaystyle\frac{\prod_{i\in S}(\lambda_w/w)_{n_i}}
                       {(\lambda_w)_n},&T\subseteq S,\\
 0,&T\not\subseteq S,
 \end{cases}
 \quad
 (u)_r\defeq\prod_{j=0}^{r-1}(u+j),\quad (u)_0=1.
\]
Here \(n_i\) counts cell \(i\), \(T=\{i:n_i>0\}\), and
\(w=|S|\). The sample is the complete ordered report, not just
an unordered count; the latter's likelihood has the extra multinomial
factor. Its next-cell conditional prediction is
\[
 P(Q_i(a_{n+1})\mid e,C_S)=
 \begin{cases}(n_i+\lambda_w/w)/(n+\lambda_w),&i\in S,\\
 0,&i\notin S.
 \end{cases}
\]
For positive total likelihood define
\[
 P(C_S\mid e)=\frac{\pi_SL_S(e)}{\sum_R\pi_RL_R(e)}.
\]
Generalization probabilities sum these posterior masses over their
constituent normal forms. Predictions average the displayed conditional
prediction over the posterior constituents.
\end{hbdefinition}

For constant \(\lambda_w=\lambda>0\), the likelihood is the
historian's equation(16), expressed without an unexplained gamma quotient:
\(\Gamma(\lambda)\prod_{i\in T}\Gamma(n_i+\lambda/w)/
[\Gamma(n+\lambda)\Gamma(\lambda/w)^c]\).
Equation(17) chooses prior weights proportional to
\(\Gamma(\alpha+w\lambda/K)/\Gamma(w\lambda/K)\).
On the present nonempty domain, \(\alpha\geq0\) gives
\[
 \pi_S=\frac{r_{|S|}}{\sum_{w=1}^K\binom Kw r_w},\qquad
 r_w=\frac{\Gamma(\alpha+w\lambda/K)}{\Gamma(w\lambda/K)}.
\]
The gamma function has its positive-real integral definition already
used in probability theory; positivity and its recurrence are also
established in the annex for these formulas. The parameter \(\alpha\)
changes constituent weights, not just the predictive pseudo-count.
The width-zero term of the source's printed sums cannot be inserted
here as a positive-probability contradictory event.

\begin{hbtheorem}{thm:p15-hintikka-countable-realization}{proved=annex:p15.hintikka-countable-realization}{Coherence and a genuinely infinite-domain realization}
The sample likelihoods are normalized, exchangeable and consistent on
deleting a last observation. They have an explicit probability-space
realization on countably many named individuals. Conditional on a
selected \(S\), every cell in \(S\) appears with probability one,
and no cell outside \(S\) appears. Consequently the event described
by \(C_S\) has probability exactly \(\pi_S\), and every
consistent monadic generalization has the sum of the corresponding
constituent masses. A proper universal exclusion can therefore have
strictly positive probability in an infinite universe.
\end{hbtheorem}

This last assertion is not achieved by replacing ``all'' with ``nearly
all''. The law \(\forall x\,Q_1x\) excludes a single contrary
individual anywhere in the named countable universe. A high frequency
of \(Q_1\) does not say that. The construction gives that exact law
a prior atom; it does not disguise a small interval near frequency one
as the universal sentence.

There is a source normalization qualification. The held history copy
calls the empty constituent contradictory on p.~319, but gives
\(1/2^K\) to all constituents in equation(18). On a nonempty domain,
the equal allocation instead normalizes to \(1/(2^K-1)\) after
excluding that impossible alternative. This is a declared coherent
reconstruction, not a claim that the printed source contained that
correction or that every original presentation uses this convention.

\section{The generalized combined system and its exact calculations}
\label{sec:p15-hintikka-combined}

\begin{hbdefinition}{def:p15-hintikka-combined-family}{Combined likelihoods and caution weights}
For the distinct generalized combined family displayed in the history's
equations(15')--(20), take \(\lambda_w=w\) and
\[
 r_w=\frac{\Gamma(\alpha+w)}{\Gamma(w)},\qquad
 \pi_S=\frac{r_{|S|}}{\sum_{u=1}^K\binom Ku r_u},
 \qquad\alpha\geq0.
\]
The conditional predictive rule is \((n_i+1)/(n+w)\) on a cell
allowed by \(S\). At \(\alpha=0\), the nonempty constituents
have equal prior weights. This is not obtained by substituting
\(\lambda=w\) indiscriminately into the different prior formula
for the constant-\(\lambda\) continuum.
\end{hbdefinition}

\begin{hbtheorem}{thm:p15-hintikka-combined-calculation}{proved=annex:p15.hintikka-combined-calculation}{Constituent, law, and predictive calculations}
For \(n\geq1\), observed support \(T\) of size \(c\), put
\[
 a_w(n,\alpha)\defeq\frac{\Gamma(\alpha+w)}{\Gamma(n+w)},
 \qquad
 D_{n,c}\defeq\sum_{w=c}^K\binom{K-c}{w-c}a_w(n,\alpha).
\]
For each \emph{particular} compatible \(S\) of width \(w\),
\[
 P(C_S\mid e)=\frac{a_w(n,\alpha)}{D_{n,c}},\qquad
 L_S(e)=\frac{(w-1)!\prod_{i\in T}n_i!}{(n+w-1)!}.
\]
The posterior of \emph{some} constituent of width \(w\) is
\(\binom{K-c}{w-c}\) times that particular-constituent probability.
For a universal law excluding a set \(B\) of \(b\) cells disjoint
from \(T\),
\[
 P(\forall x\neg\bigvee_{i\in B}Q_ix\mid e)
 =\frac{\displaystyle\sum_{w=c}^{K-b}
                  \binom{K-b-c}{w-c}a_w(n,\alpha)}{D_{n,c}}.
\]
If \(B\cap T\ne\varnothing\), that law has probability zero.
For an already observed cell \(i\in T\), prediction is
\[
 (n_i+1)\frac{\displaystyle\sum_{w=c}^K
               \binom{K-c}{w-c}a_w(n,\alpha)/(n+w)}{D_{n,c}}.
\]
For an unobserved cell \(i\notin T\), it is instead
\[
 \frac{\displaystyle\sum_{w=c+1}^K
       \binom{K-c-1}{w-c-1}a_w(n,\alpha)/(n+w)}{D_{n,c}}.
\]
Thus the observed-cell factor \(n_i+1\) cannot be blindly used for
an unobserved cell without changing the count of compatible constituents.
\end{hbtheorem}

The histogram is sufficient for the ordered likelihood; the variety
\(c\) has an additional role in eliminating constituents. The source
discusses the \emph{\(c\)-principle}, replacing the stronger demand
that a predictive representative function depend only on \(n_i,n\).
Its displayed observed-cell formula(22) must not be applied indiscriminately
to a new cell; total probability over the compatible support gives the
two different formulas above
\cite[pp.~319--325]{Niiniluoto2011HintikkaProgramReading}.

\begin{hbexample}{def:p15-hintikka-variety-case}{Same target count, different variety}
Use four cells, \(\alpha=0\), and ordered sample counts
\((2,2,0,0)\). The compatible constituents have widths 2,3,3,4.
Their unnormalized weights are \(1/120,1/360,1/360,1/840\),
so their posterior weights are \(21/38,7/38,7/38,3/38\).
The prediction for cell 1 is
\[
 3\left(\frac{21}{38\cdot6}
       +\frac{14}{38\cdot7}+\frac3{38\cdot8}\right)
 =\frac{141}{304};
\]
each unobserved cell has probability
\(7/(38\cdot7)+3/(38\cdot8)=11/304\).
With counts \((2,1,1,0)\), the compatible width 3 and width 4
constituents instead have weights \(7/10,3/10\). Cell 1 prediction
is \(3[7/(10\cdot7)+3/(10\cdot8)]=33/80\).
Both samples have \(n=4,n_1=2\), but these predictions differ.
The annex establishes these arithmetic cases as part of the calculation
theorem; they display the actual failure of the two-argument
sufficientness requirement, not a mere assertion that variety matters.
\end{hbexample}

\section{Generalization, learning, and parameter boundaries}
\label{sec:p15-hintikka-generalization}

\begin{hbtheorem}{thm:p15-hintikka-generalization-convergence}{proved=annex:p15.hintikka-generalization-convergence}{Minimal-support concentration and its precise conditions}
Suppose \(n\to\infty\) while observed support \(T\ne\varnothing\)
eventually remains fixed. With fixed positive constituent priors, both
the constant finite-\(\lambda\) family and the generalized combined
family satisfy
\[
 P(C_T\mid e)\longrightarrow1.
\]
No positive lower bound on every empirical cell frequency is required.
In the combined family the exact ratio is
\[
 \frac{a_w(n,\alpha)}{a_c(n,\alpha)}
   =\prod_{j=c}^{w-1}\frac{\alpha+j}{n+j},\qquad w>c.
\]
Every generalization entailed by \(C_T\) then has posterior tending
to one; every other generalization has posterior tending to zero.
All finite evidence compatible with two distinct constituents leaves
each of them strictly positive: this limit is not finite deductive
entailment of \(C_T\).
\end{hbtheorem}

For the combined family, a useful finite acceptance bound follows from
the same products. Let \(m=K-c\) and
\(r=(\alpha+K-1)/(n+c)\). Then
\[
 P(C_T\mid e)\geq(1+r)^{-m}.
\]
If \(m>0\), an explicit sufficient condition for posterior at least
\(1-\varepsilon\), \(0<\varepsilon<1\), is
\[
 n+c\geq\frac{\alpha+K-1}
                  {(1-\varepsilon)^{-1/m}-1}.
\]
This is a sufficient bound, not the source's exact integer threshold.
When \(c=K\), all cells have already been witnessed and the only
compatible constituent has probability one.

\begin{hbtheorem}{thm:p15-hintikka-prediction-relevance}{proved=annex:p15.hintikka-prediction-relevance}{Singular learning and instantial relevance}
For fixed finite constant \(\lambda\),
\[
 \left|P(Q_i(a_{n+1})\mid e)-n_i/n\right|\leq\lambda/n,
 \qquad n\geq1.
\]
For the combined family, the bound \((K+1)/n\) suffices.
For \(K>1\), positive priors on every nonempty constituent and
finite positive conditional parameters, another positive instance of
cell \(i\) strictly increases its next predictive probability.
The inequality is established inside this declared family; the mere
assertion of asymptotic frequency matching is not its proof.
\end{hbtheorem}

\begin{hbtheorem}{thm:p15-hintikka-caution-limits}{proved=annex:p15.hintikka-caution-limits}{Caution limits and their noncommutation}
In the combined family, increasing \(\alpha\) strictly increases
the relative prior weight of a wider constituent against a narrower one.
For fixed evidence, \(\alpha\to\infty\) concentrates the posterior
on the full-width constituent; prediction tends to
\((n_i+1)/(n+K)\), the Carnap \(c^*\) rule.
For fixed finite \(\alpha\), increasing \(n\) with \(c<K\)
instead concentrates on \(C_T\). In particular the iterated limits
of \(P(C_T\mid e)\) are
\[
 \lim_{\alpha\to\infty}\lim_{n\to\infty}P(C_T\mid e)=1,
 \qquad
 \lim_{n\to\infty}\lim_{\alpha\to\infty}P(C_T\mid e)=0.
\]
Within a fixed constituent, \(\lambda\to\infty\) in the
constant-\(\lambda\) formula gives uniform independent sampling on
its \(w\) allowed cells; \(\lambda\downarrow0\), on nonempty
evidence, gives the empirical frequency. The equal-constituent,
uniform-conditional limit realizes the Jerusalem likelihood \(w^{-n}\).
\end{hbtheorem}

The held history p.~316 reverses the straight-rule endpoint in its
prose, and gives a finite ``4'' for the no-learning endpoint despite
its displayed formula. The theorem records the actual limits of that
formula. Further, the straight empirical
guess is only a numerical limiting prediction; Reichenbach's
frequency coordination and pragmatic posit are not thereby replaced by
Hintikka's prior over universal constituents.

\section{A worked comparison of three commitments}
\label{sec:p15-hintikka-carnap-reichenbach-comparison}

\begin{hbtheorem}{thm:p15-hintikka-two-cell-comparison}{proved=annex:p15.hintikka-two-cell-comparison}{A law and an instance in the two-cell case}
Take cells \(A,\neg A\), \(n\geq1\) observed \(A\)'s, and
the combined family with \(\alpha=0\), normalized over its three
nonempty constituents. The exact universal law \(H=\forall xAx\)
and next instance have probabilities
\[
 P(H\mid e)=\frac{n+1}{n+2},\qquad
 P(A(a_{n+1})\mid e)=1-\frac1{(n+2)^2}.
\]
The equal-prior Jerusalem realization instead gives
\[
 P(H\mid e)=\frac{2^n}{2^n+1},\qquad
 P(A(a_{n+1})\mid e)=\frac{2^n+1/2}{2^n+1}.
\]
At \(n=2\), the combined values are \(3/4,15/16\); the
Jerusalem values are \(4/5,9/10\). Carnap's \(c^*\) predicts
\(3/4\) for the next instance but gives the unrestricted infinite
law probability zero. Reichenbach's straight empirical guess is one
after these two observations, but is not a prior probability of \(H\).
\end{hbtheorem}

The Carnap comparison uses the actual equal-structure construction in
Chapter~\ref{sec:p15-carnap-logical-probability-and-induction}.
The Reichenbach comparison uses the straight-rule account in
Chapter~\ref{sec:p15-reichenbach-probability-logic}.
The Hintikka law probability equals Carnap's next-instance probability
here only numerically; their targets are different. Conflating those
numbers would erase the problem that the constituent allocation solves.

\begin{hbtheorem}{thm:p15-hintikka-evidential-success}{proved=annex:p15.hintikka-evidential-success}{Concentration is not by itself truth convergence}
A nonempty world with both \(A\) and \(\neg A\) can supply an
infinite stream containing only \(A\). The displayed Hintikka law
probabilities then tend to one while \(\forall xAx\) is false.
If evidence is true and eventually witnesses every cell actually
instantiated, its eventual support is instead the true support and
the concentration theorem gives truth convergence.
If observations are independent and each of the finitely many true
cells has a fixed positive sampling probability, this witnessing
condition holds almost surely. If at least two cells are true, some
possible sample paths still omit a true cell. A finite nonempty
error-free prefix whose observed support is a proper subset of all
\(K\) cells cannot, from those reports alone, certify that no true cell
remains unseen. A prefix witnessing all \(K\) cells does certify this;
if the true support is known to be a singleton, its first true report
already witnesses the entire support.
\end{hbtheorem}

Niiniluoto calls the required truth-and-variety condition evidential
success. The epistemic calculation of a posterior and the physical
sampling assumption which makes evidence exhaustive are distinct.
One contrary report in the error-free language immediately gives
\(P(H\mid e)=0\); a hypothesis that the frequency lies near one
need not be refuted by that report. A cell never witnessed because of
selection bias is the worked failure case, not a defect repaired by
simply increasing the reported sample size
\cite[pp.~346--350]{Niiniluoto2011HintikkaProgramReading}.

\section{The axiomatic program is a further result, not just a prior choice}
\label{sec:p15-hintikka-axiomatic-program}

Niiniluoto's historical account states six distinct structural demands.
Probability axioms constrain degrees as probabilities. Finite regularity
excludes certainty in a contingent singular conclusion not entailed by
its singular evidence. Individual symmetry permits permutations of names;
cell symmetry permits permutations of the exhaustive \(Q\)-cells.
The stronger sufficientness or \(\lambda\)-principle requires a
representative function \(f(n_i,n)\), independent of how the remaining
count is distributed among other cells. The weaker \(c\)-principle
allows \(f(n_i,n,c)\). Neither symmetry requirement alone entails
either sufficientness condition. The four-cell calculation exhibits that
difference concretely, while preserving both symmetries.

The history reports a Kemeny--Carnap characterization under its
probability, regularity, symmetry and two-argument conditions; its
two-cell version needs an additional linearity hypothesis. It then
reports the Hintikka--Niiniluoto \(K\)-dimensional characterization,
with a parameter controlling full-width conditional prediction and
separate \(\gamma_c=f(0,c,c)\) controlling the probability of
finding a new cell after \(c\) distinct observations. These parameters
are not interchangeable with the single \(\alpha\) of the earlier
continuum. The reported result explains how a predictive sufficientness
assumption can force a skeptical prior on generalizations; it is not
established merely by supplying the sufficient example constructed here
\cite[secs.~3--4, pp.~322--326]{Niiniluoto2011HintikkaProgramReading}.

In the history's notation the conditional rule of the later family is
\[
 P(Q_i(a_{n+1})\mid e,C_S)
       =\frac{n_i+\lambda/K}{n+w\lambda/K},\qquad i\in S.
\]
This corresponds to the \emph{declared} conditional construction here
with \(\lambda_w=w\lambda/K\), when \(\lambda>0\).
The source's full stated parameter range, necessity argument and
constituent-prior restrictions are not supplied by that substitution.
Its \(\gamma_c\) upper bound and the ratio formula for the
atomistic prior belong to the full characterization, not to a theorem
that every freely chosen vector of \(\gamma_c\)'s is realizable.
The original 1976 result and its 1980 postscript remain distinct proof
and version dependencies.

The history also distinguishes equal weight per constituent from equal
weight per \emph{constituent-structure}, the sentence that exactly
\(w\) cells are populated. Giving widths equal mass would give each
particular width-\(w\) constituent mass proportional to
\(1/\binom Kw\), not equal mass to every constituent. This is a
different inductive commitment, even though it still treats individual
cell names symmetrically. The wider axiomatic program is meant to
describe such choices rather than identify every symmetry with the
earliest equal-constituent allocation.

\section{Information, confirmation and the acceptance of a generalization}
\label{sec:p15-hintikka-information-acceptance}

The history connects Hintikka's program with a distinction between high
posterior probability and positive relevance. A law may be improbable
before an observation and yet receive an increase in probability from
it. Conversely a previously very probable hypothesis may remain highly
probable while being disconfirmed. Information and acceptance have
further typed roles; their names do not turn them into the same degree.

\begin{hbdefinition}{def:p15-hintikka-information-functions}{Information and corroboration on a fixed constituent space}
For a hypothesis \(h\) with \(P(h)>0\), define surprise
\(\operatorname{inf}(h)=-\log P(h)\); substantial content is
\(\operatorname{cont}(h)=1-P(h)\). For defined positive
conditionals, conditional surprise is \(-\log P(h\mid e)\),
and transmitted surprise is
\(\log[P(h\mid e)/P(h)]\). Transmitted conditional content is
\(P(h\mid e)-P(h)\). Additive surprise of \(h\) beyond \(e\) is
\(\operatorname{inf}(h\land e)-\operatorname{inf}(e)\), where
both logarithms are defined. Additive content is instead
\(\operatorname{cont}(h\land e)-\operatorname{cont}(e)\).
These are different differences, not two names for one quantity.

For a nonfalse generalization with constituent normal form
\(h=\bigvee_{S\in I_h}C_S\), the history's minimum
corroboration function is
\[
 \operatorname{corr}(h,e)\defeq
               \min_{S\in I_h}P(C_S\mid e).
\]
Its range includes zero when \(h\)'s normal form includes a constituent
refuted by \(e\). It is not the sum which defines \(P(h\mid e)\).
\end{hbdefinition}

\begin{hbtheorem}{thm:p15-hintikka-information-laws}{proved=annex:p15.hintikka-information-laws}{The different orderings and their countercases}
Where the logarithms are defined, additive surprise equals conditional
surprise, and transmitted surprise is positive exactly for positive
relevance. For global logical entailment \(h_1\models h_2\),
\[
 P(h_1\mid e)\leq P(h_2\mid e),\qquad
 \operatorname{corr}(h_1,e)\geq\operatorname{corr}(h_2,e).
\]
Under minimal-support concentration, \(C_T\) eventually maximizes
this minimum-corroboration function among all nonfalse generalizations.
The corroboration inequality need not hold if global entailment is
replaced by entailment \emph{only relative to evidence} without changing
the definition's constituent index set.
\end{hbtheorem}

That last condition preserves an actual qualification in using the
history's displayed equation(42): its evidence-relative formulation
does not follow for a minimum over all original disjuncts. An excluded
disjunct has zero posterior and remains in that minimum. Conditioning
away the excluded disjunct would be a changed definition and must be
stated, not attributed silently to the formula.

The history presents acceptance as a further methodological choice:
after a sufficiently large sample, accept the minimal constituent whose
posterior exceeds a chosen threshold, and accept its deductive
consequences. This avoids accepting many unrelated high-probability
sentences independently, the pattern responsible for lottery-style
inconsistency. Its proposed reduction of singular acceptance to
acceptance of a generalization is likewise a rule of policy, not a
deductive consequence of one high predictive value. The finite bound
above can implement a sufficient threshold on the declared family;
it does not certify that the evidence is exhaustive or error free
\cite[secs.~5--7, pp.~329--337]{Niiniluoto2011HintikkaProgramReading}.

\begin{hbtheorem}{thm:p15-hintikka-cognitive-utility}{proved=annex:p15.hintikka-cognitive-utility}{Expected truth, content-sensitive payoff, and a source qualification}
On a finite exhaustive constituent space, let accepting \(g\) give
payoff 1 when true and 0 when false. Its expected payoff on \(e\)
is \(P(g\mid e)\), so a tautology always maximizes it.
If instead the payoff is \(1-P(g)\) when \(g\) is true and
\(-P(g)\) when it is false, its expected payoff is
\(P(g\mid e)-P(g)\).
If the false payoff is changed to \(-qP(g)\),
\(0<q\leq1\), while the true payoff remains \(1-P(g)\),
the expectation is
\[
 P(g\mid e)[1-(1-q)P(g)]-qP(g),
\]
not in general \(P(g\mid e)-qP(g)\).
\end{hbtheorem}

For example \(P(g)=1/4\), \(P(g\mid e)=3/4\),
\(q=1/2\) gives the actual expectation \(17/32\), whereas
the shortened expression gives \(5/8\). The history's equations(54)
and(55) have this mismatch; they agree at \(q=1\).
An alternative payoff with \(1-qP(g)\) for truth would produce
the shortened expression, but that would be a different specified
utility. The utility theorem proves arithmetic conditional on the
payoffs; it does not prove a unique epistemic value system or vindicate
the choice of its boldness weight.

\section{The program's extensions have their own assumptions}
\label{sec:p15-hintikka-extensions}

The history's identity extension replaces a simple populated/unpopulated
cell by cardinality alternatives \(0,1,\ldots,d-1,\text{at least }d\)
at depth \(d\). Such descriptions split further at the next depth;
they are not the same objects as the unary constituents without
identity. Its ordered-universe and relational developments likewise
require information about successors or pairs, rather than only one
histogram. The original normal-form construction explains how such
information can be expressed; it does not automatically prove the
probabilistic adequacy or the general consistency filtering needed for
all those extensions.

For nomic or causal modalities the history describes Uchii's S5-style
constituents, recording which profiles are physically possible and
excluding the rest necessarily. Actual existence implies possibility,
but the converse is an additional principle. The discussion of
plenitude objects precisely to obtaining certainty about nomic
possibility from certainty about the actual support by a premise which
already makes every relevant possibility realized. No result above
licenses that move from extensional description to causal law
\cite[sec.~4, pp.~326--329]{Niiniluoto2011HintikkaProgramReading}.

Theoretical support is distinct from observational support. The history
considers \(P(h\mid e,T)\), with a background theory \(T\),
and distinguishes comparing it with \(P(h)\) from comparing it with
\(P(h\mid T)\). These are different baselines. Its critique of
Hempel's theoretical-predicate example concerns combining converse
entailment and special-consequence principles of confirmation as though
one notion satisfied both. The reported indispensability result for
theoretical terms needs the 1973 Niiniluoto--Tuomela constructions;
it is not proven by adding an extra condition to Bayes's notation.
Similarly, changing \(\alpha\) or \(\lambda\) changes the
probability model, and is not ordinary conditionalization within the
old fixed model unless a hierarchical parameter construction is
separately specified
\cite[sec.~8, pp.~337--340]{Niiniluoto2011HintikkaProgramReading}.

The analogy discussion distinguishes exact same-cell repetition from
similarity between different profiles. A Hamming fraction over the
primitive predicates supplies one possible distance, and families of
attributes may have their own metrics. The historical proposals for
analogy profit add evidence-dependent terms to predictive counts;
some earlier multipliers failed instantial relevance or asymptotic
frequency matching. Such modifications must establish normalization,
coherence and the desired limiting conditions anew. The basic symmetric
construction here has no hidden metric and does not solve negative
analogy merely by calling two cells similar.

Observational error also changes the likelihood. The history
distinguishes \(Q_j(a)\), the true cell, from \(S_i(a)\), its
reported appearance, and introduces \(P(S_i(a)\mid Q_j(a))\).
The error-free sample likelihood here assumes reports are true; one
apparent contrary report cannot refute a true constituent with the
same certainty if misclassification is allowed. A declared error matrix
and a sample-error dependence model are needed for that extension
\cite[sec.~9, pp.~340--343]{Niiniluoto2011HintikkaProgramReading}.

\section{Reports with a declared error channel}
\label{sec:p15-hintikka-error-channel}

The distinction between true cells and reported appearances can be made
into a complete probability construction. The following extension retains
the constituent prior and conditional sample law already declared. It adds
a fixed, explicit observation mechanism; it is a reconstruction of the
error problem, not a claim to reproduce every historical error calculus
\cite[sec.~9, pp.~340--343]{Niiniluoto2011HintikkaProgramReading}.

\begin{hbdefinition}{def:p15-hintikka-error-channel}{True cells and noisy reports}
Let the report alphabet be \(\{1,\ldots,J\}\), where \(J\geq1\).
Specify numbers \(E_{ri}\geq0\) with
\(\sum_{r=1}^J E_{ri}=1\) for every true cell \(i\).
After the latent sequence of true cells has been sampled, reports are
conditionally independent, and a true cell \(i\) produces report \(r\)
with probability \(E_{ri}\). The same matrix is used at every position.
This is an assumption about observation, separate from the conditional
dependence among true cells in the constituent sample law.

For a support \(S\), write \(w=|S|\),
\(a_S=\lambda_w/w\), and let \(\boldsymbol c=(c_i)_{i\in S}\)
be latent cell counts. For the first \(n\) actual reports
\(\boldsymbol y=(y_1,\ldots,y_n)\), define weights
\(W_{S,n}(\boldsymbol c;\boldsymbol y)\) by
\(W_{S,0}(\boldsymbol0;\varnothing)=1\). All other initial weights
are zero. Given the weights at stage \(n\), start the next array at
zero and, for every \(i\in S\), add
\[
 W_{S,n}(\boldsymbol c;\boldsymbol y)
 \frac{c_i+a_S}{n+\lambda_w}E_{y_{n+1},i}
\]
to \(W_{S,n+1}(\boldsymbol c+\boldsymbol e_i;
\boldsymbol y,y_{n+1})\). Here \(\boldsymbol e_i\) is the
unit count vector; it is not evidence. Only vectors with total count
\(n\) occur at stage \(n\), so each array is finite.
\end{hbdefinition}

\begin{hbtheorem}{thm:p15-hintikka-error-channel}{proved=annex:p15.hintikka-error-channel}{Coherent noisy reports and exact finite updating}
The true-cell construction and the declared report channel have a joint
countable-domain realization with unchanged constituent masses \(\pi_S\).
The report likelihood conditional on \(C_S\) is
\[
 H_S(\boldsymbol y)=\sum_{\boldsymbol c}W_{S,n}
       (\boldsymbol c;\boldsymbol y)
 =\sum_{(i_1,\ldots,i_n)\in S^n}
       L_S(i_1,\ldots,i_n)\prod_{t=1}^n E_{y_t,i_t}.
\]
These likelihoods are normalized, deletion consistent and exchangeable.
For \(Z(\boldsymbol y)=\sum_S\pi_SH_S(\boldsymbol y)>0\),
the joint posterior of support and latent count vector is
\[
 v_{S,\boldsymbol c}
   =\frac{\pi_SW_{S,n}(\boldsymbol c;\boldsymbol y)}
                {Z(\boldsymbol y)}.
\]
The next true-cell prediction is
\(\sum_{S\ni i,\boldsymbol c}v_{S,\boldsymbol c}
(c_i+a_S)/(n+\lambda_{|S|})\); the next report prediction is
\[
 P(Y_{n+1}=r\mid\boldsymbol y)
  =\sum_{S,\boldsymbol c}v_{S,\boldsymbol c}
       \sum_{i\in S}\frac{c_i+a_S}{n+\lambda_{|S|}}E_{ri}.
\]
Every constituent has positive posterior after every finite report string
if every \(E_{ri}>0\). If all columns of the channel are the same
probability vector \((q_r)\), reports leave the constituent priors
unchanged whenever their likelihood is positive. Thus coherence of error
correction does not imply that the reports identify the true support.
\end{hbtheorem}

\begin{hbexample}{def:p15-hintikka-noisy-counterreport}{One apparent counterexample and complete loss of information}
Take the three nonempty supports on \(A,\neg A\), equal prior masses,
and the combined family. Let a report flip its true cell with probability
\(0<\varepsilon<1/2\). After one report of \(\neg A\), the
conditional likelihoods on the \(A\)-only, \(\neg A\)-only and
full supports are respectively \(\varepsilon,1-\varepsilon,1/2\).
Consequently
\[
 P(\forall xAx\mid Y_1=\neg A)=\frac{2\varepsilon}{3}.
\]
At \(\varepsilon=1/10\) this is \(1/15\), rather than zero.
The report weakens the law without deductively refuting it. At
\(\varepsilon=1/2\), both channel columns are \((1/2,1/2)\);
every length-\(n\) report string has likelihood \(2^{-n}\) under
every support. No number of such reports changes any support probability.
These conclusions concern the specified channel; an unknown or dependent
error process requires a different model and a new likelihood.
\end{hbexample}


\section{Truthlikeness and the final philosophical assessment}
\label{sec:p15-hintikka-truthlikeness}

\begin{hbdefinition}{def:p15-hintikka-truthlikeness-space}{A distinct target and distance structure}
Let \(B=\{h_1,\ldots,h_m\}\), \(m\geq2\), be finite,
exhaustive and exclusive complete answers, and let
\(\Delta_{ij}\in[0,1]\) be zero exactly when \(i=j\).
For a nonempty index set \(I_g\subseteq\{1,\ldots,m\}\) and
partial answer \(g=\bigvee_{j\in I_g}h_j\),
set
\[
 \Delta_{\min}(h_i,g)=\min_{j\in I_g}\Delta_{ij},\qquad
 \Delta_{\mathrm{sum}}(h_i,g)=
              \frac{\sum_{j\in I_g}\Delta_{ij}}{\sum_{j=1}^{m}\Delta_{ij}}.
\]
The denominator is positive. With chosen positive weights \(s,t\),
put \(\operatorname{Tr}(g,h_i)=1-s\Delta_{\min}-t\Delta_{\mathrm{sum}}\).
Its expected degree on evidence is
\(\operatorname{ver}(g\mid e)=\sum_iP(h_i\mid e)\operatorname{Tr}(g,h_i)\).
The unknown true answer is the target; the posterior expectation is
an estimate of that truthlikeness, not its actual value.
\end{hbdefinition}

\begin{hbtheorem}{thm:p15-hintikka-truthlikeness-limit}{proved=annex:p15.hintikka-truthlikeness-limit}{Expected truthlikeness and probable approximate truth}
The displayed truthlikeness equals one exactly for the complete target
answer. If posterior mass concentrates on \(h_*\), then
\(\operatorname{ver}(g\mid e)\to\operatorname{Tr}(g,h_*)\).
For \(\varepsilon\geq0\), define the blurred answer by adding
every \(h_i\) at minimum distance at most \(\varepsilon\)
from \(g\). Its probability, or probable approximate truth, is
\[
 \operatorname{PAT}_{1-\varepsilon}(g\mid e)
   =\sum_{i:\Delta_{\min}(h_i,g)\leq\varepsilon}P(h_i\mid e)
   \geq P(g\mid e).
\]
For a finite space it tends to \(P(g\mid e)\) as
\(\varepsilon\downarrow0\). A refuted answer may have positive
probable approximate truth and positive expected truthlikeness.
These are different properties from positive probability of literal truth.
\end{hbtheorem}

In a two-answer space at mutual distance \(1/4\), let evidence
make \(h_2\) certain and refute \(g=h_1\). With \(s=t=1/4\),
its actual and expected truthlikeness is \(11/16\), while its
posterior truth probability is zero. At tolerance \(1/4\), its
probable approximate truth is one. The example makes the distinction
usable; it does not infer the appropriate distance or utility weights
from the probabilities themselves
\cite[secs.~10--11, pp.~343--348]{Niiniluoto2011HintikkaProgramReading}.

The historian's concluding argument treats the succession of systems
as a progressive methodological research program: later systems address
problems that the earlier ones leave aside, notably genuine
generalization. He nevertheless separates this assessment from a claim
that every theory's numerical degree is mechanically available. The
program's later interrogative model treats Nature's answers as possibly
general or quantificational, and questions restricting inquiry to
atomic observation reports. Niiniluoto doubts that explicitly formulated
regularity assumptions remove every inductive step. A premise which
already guarantees the requisite regularity does not acquire an
independent justification by appearing in a deductive argument
\cite[sec.~12, pp.~349--352]{Niiniluoto2011HintikkaProgramReading}.

\begin{hbmethod}{def:p15-hintikka-reasoning-procedure}{Using the construction without erasing its commitments}
Choose the language, domain convention and classification; identify
the exact law or singular prediction sought. Form the constituent normal
forms before assigning masses. Declare the constituent weights and
conditional predictive family, distinguishing constant-\(\lambda\)
from width-dependent combined rules. Record the complete sample counts
and observed variety; eliminate supports contradicted by error-free
evidence. Compute likelihoods, normalize posterior constituent masses,
then sum for laws or average for predictions. Keep acceptance thresholds,
epistemic utilities, distances and physical sampling assumptions as
additional commitments. For erroneous reports with a known fixed channel, use the joint latent-count
recursion in \cref{thm:p15-hintikka-error-channel}. For a different error
process, relations or identity in the language, or a changed conceptual
scheme, supply the corresponding construction instead of transferring
the error-free formulas unchanged.
\end{hbmethod}

\begin{proofboundary}{def:p15-hintikka-source-boundary}{Declared constructions and wider characterization problems}
The annex proves the stated normal-form constructions, explicit countable
probability realization, likelihood/posterior computations, convergence
and its countercases, parameter limits, and displayed information,
utility and distance calculations. These are complete proofs on their
declared domains, not proofs of the full historical program.
Normal-form existence alone gives neither probability weights nor a
satisfiability filter for unrestricted relational descriptions. The
K-dimensional necessity/representation theorem and identity, relational,
modal, analogy and theoretical-systematization extensions require
their own hypotheses and arguments. They do not follow by substituting
new parameter names into the monadic likelihood. The fixed observation-channel extension now has its own joint realization,
likelihood, updating proof and nonidentification counterexample; it does not
cover an unspecified or dependent error process. The historical author copy
and final 2011 publisher version retain their distinct identities.
\end{proofboundary}
